\documentclass[10pt]{article}
\usepackage[T1]{fontenc}
\usepackage{lmodern,amsmath,amssymb,amsthm,microtype}
\usepackage[margin=.8in]{geometry}
\usepackage[colorlinks=true,urlcolor=blue]{hyperref}
\newtheorem{theorem}{Theorem}
\title{The Exponential Formula and Three Functional Identities of A092635}
\author{Qichao Wang\\{\small Hebei University of Technology}\\[.3em]Daoyu Dong\\{\small University of Electronic Science and Technology of China}}
\date{October 26, 2026}
\begin{document}\maketitle
\begin{abstract}
We prove the Euler-polynomial exponential formula proposed for
A092635 and all three functional identities conjectured in its
2024 update. The arguments use the formal asymptotic expansion
of the gamma ratio defining the sequence, oddness of its logarithm,
and the gamma recurrence. The identities hold to every order;
no convergence of the generally divergent generating series is
assumed.
\end{abstract}
\section{Definition and exponential formula}
OEIS A092635~\cite{oeis} defines the coefficients $a_j$ by
\[
 F(t)=\frac{\pi t}{16^t}\binom{2t}{t}^{\!2}
 \sim G\bigl(1/(8t)\bigr),\qquad
 G(x)=\sum_{j\ge0}a_jx^j,\quad a_0=1.
\]
All series identities below are formal. Define the integers $b_j$ by
\[
 B(z)=\frac2{1+e^{4z}}=\sum_{j\ge0}b_j\frac{z^j}{j!}.
\]
The entry proposes $G(x)=\exp(\sum_{j\ge1}b_jx^j/j)$ and its
coefficient recurrence; it separately conjectures three identities
for $H(x)=xG(x)$.
\begin{theorem}
In $\mathbb Q[[x]]$,
\[
 G(x)=\exp\left(\sum_{j\ge1}\frac{b_j}{j}x^j\right),
 \qquad
 a_m=\frac1m\sum_{j=1}^m b_j a_{m-j}\quad(m\ge1).
\]
Here $b_j=(-4)^jE_j(1)$, with Euler polynomials defined by
$2e^{uz}/(e^z+1)=\sum E_j(u)z^j/j!$.
\end{theorem}
\begin{proof}
Let $B_{2k}$ be the Bernoulli numbers in Stirling's logarithmic
expansion~\cite{dlmf},
\[
 \log\Gamma(t+1)\sim
 (t+\tfrac12)\log t-t+\tfrac12\log(2\pi)
 +\sum_{k\ge1}\frac{B_{2k}}{2k(2k-1)t^{2k-1}}.
\]
Substitution in $\log F(t)$ cancels the elementary terms and gives
\[
 \log G(x)=\sum_{k\ge1}
 \frac{2B_{2k}(2^{-(2k-1)}-2)8^{2k-1}}
 {2k(2k-1)}x^{2k-1}.
\]
On the other hand, use
$2/(1+e^w)=2/(e^w-1)-4/(e^{2w}-1)$ and the Bernoulli
generating function. For $j\ge1$ it gives
\[
 b_j=\frac{2(1-2^{j+1})4^jB_{j+1}}{j+1}.
\]
This is zero for positive even $j$, and its odd coefficients
match $j[x^j]\log G(x)$ exactly. Hence the proposed exponential
formula holds. Formal logarithmic differentiation gives the
coefficient recurrence. Finally
$B(z)=2e^{-4z}/(1+e^{-4z})$ is the defining Euler-polynomial
generating function at $u=1$, with argument $-4z$.
\end{proof}

\section{The three reciprocal and substitution identities}
\begin{theorem}
For $H(x)=xG(x)$ the following three conjectures hold as formal
Laurent-series identities:
\[
 H(x)=-\frac{x^2}{H(-x)},\qquad
 H(x)=-H\left(\frac{-x}{1+4x}\right),\qquad
 H(x)=\frac{x^2}{H(x/(1-4x))}.
\]
\end{theorem}
\begin{proof}
The logarithm derived above contains only odd powers. Thus
$G(x)G(-x)=1$, or $H(x)H(-x)=-x^2$, proving the first identity.
This is a formal conclusion about coefficients; it does not
evaluate the analytic gamma expression at negative $t$.

The gamma duplication formula gives, for positive $t$,
\[
 F(t)=t\left(\frac{\Gamma(t+1/2)}{\Gamma(t+1)}\right)^2.
\]
The recurrence $\Gamma(t+1)=t\Gamma(t)$ therefore yields the
exact identity
\[
 F(t)F(t-1/2)=\frac{t-1/2}{t}.
\]
Asymptotic expansions at positive infinity are unique and may be
multiplied and shifted by a fixed constant. With $x=1/(8t)$ this
becomes the formal identity
\[
 G(x)G\left(\frac{x}{1-4x}\right)=1-4x.
\]
Multiplying by $x^2/(1-4x)$ proves
$H(x)H(x/(1-4x))=x^2$, the third assertion.
Combine it with the first to obtain
$H(x/(1-4x))=-H(-x)$. Replacing $x$ by $-x$ proves the second
assertion. All substituted series have zero constant term, so
their formal compositions are defined.
\end{proof}
\paragraph{Verification and scope.}
The standard-library verifier computes the Euler coefficients and
the proposed exponential recurrence exactly with rational arithmetic,
and compares all 301 published coefficients. Independently, the
exact gamma recurrence
$F(t+1)/F(t)=(t+1/2)^2/(t(t+1))$ determines each next asymptotic
coefficient by coefficient comparison; the verifier reconstructs
all 301 terms that way as a separate control. It verifies all
three identities through degree 100 using exact formal operations.
The proof establishes the identities at every order. We do not
claim convergence, new authorship of Stirling's formula, or OEIS
acceptance of an update.
\begin{thebibliography}{2}
\bibitem{oeis} M. Somos, OEIS A092635, including the proposals of
P. Bala (2015) and P. D. Hanna (2024),
\url{https://oeis.org/A092635}, accessed October 26, 2026.
\bibitem{dlmf} NIST Digital Library of Mathematical Functions,
equation 5.11.1, logarithmic Stirling expansion,
\url{https://dlmf.nist.gov/5.11.E1}.
\end{thebibliography}\end{document}
